\section{NVMe 同盘 repo（E:）十重复现}\label{sec:nvme-e}

\begin{table}[htbp]
\centering
\caption{E: repo 十次墙钟（2026-07-10，warm-up 后）}
\label{tab:e-raw}
\small
\begin{tabular}{@{}r rr@{}}
\toprule
Run & $T_i$ (s) & $v_i$ (MB/s) \\
\midrule
 1 & 19.494 & 263.08 \\
 2 & 18.919 & 271.08 \\
 3 & 18.909 & 271.22 \\
 4 & 18.726 & 273.87 \\
 5 & 19.357 & 264.95 \\
 6 & 18.941 & 270.76 \\
 7 & 18.817 & 272.54 \\
 8 & 19.091 & 268.63 \\
 9 & 19.038 & 269.38 \\
10 & 18.961 & 270.47 \\
\midrule
均值 $\bar{T}$ & 19.025 & 269.56 \\
样本标准差 $s$ & 0.236 & — \\
$\mathrm{CV}$ & \multicolumn{2}{c}{1.24\%} \\
极差 & 18.726 -- 19.494 & 263.08 -- 273.87 \\
\bottomrule
\end{tabular}
\end{table}

\begin{figure}[htbp]
\centering
\begin{tikzpicture}
\begin{axis}[
  width=0.92\linewidth,
  height=6cm,
  ybar,
  bar width=14pt,
  ymin=260, ymax=276,
  ylabel={$v_i$ (MB/s)},
  xlabel={Run},
  xtick=data,
  symbolic x coords={1,2,3,4,5,6,7,8,9,10},
  nodes near coords,
  nodes near coords align={vertical},
  every node near coord/.append style={font=\scriptsize},
  grid=major,
  grid style={dashed,gray!30},
  legend style={at={(0.5,1.12)},anchor=south,legend columns=2,font=\footnotesize}
]
\addplot[fill=NVMeColor!35,draw=NVMeColor] coordinates {
  (1,263.08) (2,271.08) (3,271.22) (4,273.87) (5,264.95)
  (6,270.76) (7,272.54) (8,268.63) (9,269.38) (10,270.47)
};
\addlegendentry{E: NVMe repo}
\draw[dashed,SatRed,thick] (axis cs:1,269.56) -- (axis cs:10,269.56);
\addlegendimage{dashed,SatRed,thick}\addlegendentry{$\bar{v}=269.56$}
\end{axis}
\end{tikzpicture}
\caption{E: 十次有效吞吐 $v_i$；虚线为均值。}
\label{fig:e-bar}
\end{figure}

\begin{lemma}[E: 低方差稳态]
\label{lem:e-low-cv}
在 \cref{asm:steady} 下，E: repo 满足 $\mathrm{CV}=1.24\% < 3\%$，且
$v_i\in[263.08,\,273.87]$，带宽 $\Delta v/\bar{v}=4.00\%$。
\end{lemma}

\begin{proof}
由 \cref{tab:e-raw} 直接计算 \cref{eq:cv}。
\end{proof}

\section{NVMe 同盘异分区对照（D:/C:）}\label{sec:partition}

\begin{table}[htbp]
\centering
\caption{同盘异分区 repo 十次汇总（源均为 D:）}
\label{tab:partition}
\small
\begin{tabular}{@{}l rrr rr@{}}
\toprule
Repo & $\bar{T}$ (s) & $s$ (s) & CV (\%) & $\bar{v}$ (MB/s) & $v$ 极差 \\
\midrule
E: 数据卷 & 19.025 & 0.236 & 1.24 & 269.56 & 4.00\% \\
D: 源同分区 & 18.987 & 1.349 & 7.11 & 270.11 & 3.4\% \\
C: 系统卷 & 19.616 & 1.193 & 6.08 & 261.44 & 4.8\% \\
\bottomrule
\end{tabular}
\end{table}

\begin{table}[htbp]
\centering
\caption{D: repo 原始墙钟（秒）}
\label{tab:d-raw}
\footnotesize
\begin{tabular}{@{}*{10}{c}@{}}
\toprule
1 & 2 & 3 & 4 & 5 & 6 & 7 & 8 & 9 & 10 \\
\midrule
18.292 & 18.776 & 18.378 & 22.749 & 18.496 & 18.434 & 18.223 & 18.640 & 19.137 & 18.742 \\
\bottomrule
\end{tabular}
\end{table}

\begin{table}[htbp]
\centering
\caption{C: repo 原始墙钟（秒）}
\label{tab:c-raw}
\footnotesize
\begin{tabular}{@{}*{10}{c}@{}}
\toprule
1 & 2 & 3 & 4 & 5 & 6 & 7 & 8 & 9 & 10 \\
\midrule
18.942 & 19.592 & 19.348 & 18.615 & 19.128 & 22.874 & 18.982 & 19.618 & 19.503 & 19.561 \\
\bottomrule
\end{tabular}
\end{table}

\begin{figure}[htbp]
\centering
\begin{tikzpicture}
\begin{axis}[
  width=0.88\linewidth,
  height=6.2cm,
  ybar,
  bar width=22pt,
  ymin=255, ymax=275,
  ylabel={$\bar{v}$ (MB/s)},
  symbolic x coords={E:,D:,C:},
  xtick=data,
  nodes near coords,
  grid=major,
  grid style={dashed,gray!30},
  enlarge x limits=0.25
]
\addplot[fill=NVMeColor!40,draw=NVMeColor]
  coordinates {(E:,269.56) (D:,270.11) (C:,261.44)};
\end{axis}
\end{tikzpicture}
\caption{同盘异分区均值吞吐（E:/D:/C:）。}
\label{fig:partition-bar}
\end{figure}

\begin{proposition}[分区弱不变性（同物理 NVMe）]
\label{prop:partition}
在 Disk 0 上，E:/D:/C: 三处 repo 的 $\bar{v}$ 落在
$261.44\text{--}270.11$\,MB/s，相对 spread
$\frac{270.11-261.44}{269.56}\approx 3.2\%$。
\end{proposition}

\begin{proof}
\cref{tab:partition} 直接比较；均值差异小于单次 run 的典型 jitter。
离群点 Run4（D: 22.749\,s）与 Run6（C: 22.874\,s）归因 OS/调度噪声，不改变 $\bar{v}$ 同阶结论。
\end{proof}
